% Generated draft: requires mathematical and source review.
\documentclass[11pt]{article}
\usepackage[margin=25mm]{geometry}
\usepackage{fontspec}
\setmainfont{Latin Modern Roman}
\usepackage{amsmath,amssymb,mathtools,unicode-math}
\setmathfont{Latin Modern Math}
\usepackage{enumitem,needspace,xurl,hyperref}
\setlength{\parindent}{0pt}
\setlength{\parskip}{5pt}
\setlength{\emergencystretch}{3em}
\newcommand{\N}{\mathbb{N}}
\newcommand{\Z}{\mathbb{Z}}
\newcommand{\Q}{\mathbb{Q}}
\newcommand{\R}{\mathbb{R}}
\newcommand{\C}{\mathbb{C}}
\newcommand{\F}{\mathbb{F}}
\newcommand{\A}{\mathbb{A}}
\newcommand{\PP}{\mathbb{P}}
\newcommand{\E}{\mathbb{E}}
\newcommand{\eps}{\varepsilon}
\newcommand{\dd}{\,\mathrm{d}}
\begin{document}
\section*{UnsolvedMath ID: 801067}
\textbf{Source problem code:} catalogue entry
\subsection{Definitions and mathematical statement}
Let \(f\) be a H\'enon map in the setting intended by the source, and let \(X\) denote its phase space. The supplied record does not specify the ambient domain, the allowed class of maps, or the definition of quasi-hyperbolicity; these remain unresolved qualifications of the statement below. Boundedness, interior, boundary, and attraction are understood in the intended ambient setting.

Define
\[
K^\pm=\{z\in X:\{f^{\pm n}(z):n\geq 0\}\text{ is bounded}\},\qquad J^\pm=\partial K^\pm,\qquad J=J^+\cap J^-.
\]
The source also defines \(J^*=\overline{\mathcal S}\) and states that \(J^*\subset J\), but does not define \(\mathcal S\) in the supplied record.

A sink orbit is an attracting periodic orbit of \(f\). For such an orbit \(\mathcal O\), write
\[
B_f(\mathcal O)=\{z\in X:\operatorname{dist}(f^n(z),\mathcal O)\longrightarrow 0\text{ as }n\longrightarrow\infty\},
\]
where distance refers to the intended ambient metric.

The question is: if \(f\) is quasi-hyperbolic, must there exist a finite family of sink orbits \(\mathcal O_1,\ldots,\mathcal O_r\), with \(r\geq 0\), such that
\[
\operatorname{int}_X(K^+)=\bigcup_{j=1}^{r}B_f(\mathcal O_j)?
\]
The empty family is permitted when \(r=0\). The hypothesis of quasi-hyperbolicity is retained verbatim and requires a source-supported definition before this draft is a fully specified mathematical statement.
\subsection{Short English statement}
In the source's intended H\'enon-map setting, does quasi-hyperbolicity imply that the interior of the forward bounded-orbit set consists of finitely many basins of attracting periodic orbits? The ambient setting and the meaning of quasi-hyperbolicity require clarification.
\subsection{Sources}
\begin{itemize}
\item Local UnsolvedMath registry, permanent ID \(801067\), exact source problem code \texttt{catalogue entry}, titled ``Henon Maps: Open Problems --- Question 19.'' The record attributes the question to Berger and Bedford (2023) and labels it open; this is a supplied status label, not an independently verified mathematical assessment.
\item Berger and Bedford, \emph{Henon Maps: A List of Open Problems} (2023), Question 19, \url{https://arxiv.org/abs/2312.03907}. The supplied provenance records access on 2026-10-16 and extraction by \texttt{arxiv\_latex}; these provenance details have not been independently verified here.
\item The statement cites reference [BS] for standard definitions. The supplied record provides no bibliographic details identifying [BS].
\end{itemize}
\end{document}
